%% ============================================================
%%  §6 代价、局限与威胁效度
%% ============================================================

\section{代价、局限与威胁效度}
\label{sec:limitations}

本节做两件事：先说清这套设计的\textbf{代价}（属于 MLIR，不属于本文方法），再说清本文证据
\textbf{能支撑到哪一步}，并把结论按强度分级。

\subsection{机制层面的代价}

\paragraph{栈顶仍然碎片化}
第~\ref{sec:evidence} 节的投入分布显示，导入层只占已归属文件的 4.5\%，而承接框架的工作分散在
\texttt{torch-mlir}、StableHLO、\texttt{onnx-mlir} 等\textbf{不同仓库}里。共享的收益并没有在栈顶
兑现，而是被推迟到了各生态自己的项目中。准确的说法不是"MLIR 没有解决框架接入"，而是
"\textbf{MLIR 把栈顶的共享收益外置给了生态}"。

\paragraph{抽象泄漏集中在三处}
层与层之间的边界并非处处等价。最难的几处是：\textbf{bufferization}~\cite{mlirbufferization}（从值语义到缓冲区语义，
每次读写都要决定 copy 还是原地改写，这是语义决策而非语法改写）、\textbf{shape 推断}~\cite{mlirshape}
（动态维度需要约束与见证的传播，语义横跨 shape、tensor、memref 三层）、\textbf{量化}
（scale 与 zero-point 的归属需要额外的类型层承载）。这类信息无法被完全抽象掉，
因此每增加一层就增加一处潜在的语义错位点。

\paragraph{pass 链变长，调试面变大}
第~\ref{sec:mechanisms} 节机制四指出边界是可声明的，代价是 pass 顺序变敏感、不完整的下降会留下
难以定位的残留。需要明确：\textbf{本文没有测量这一成本}，它只作为机制层面的定性判断存在。

\paragraph{治理成本}
第~\ref{sec:contract} 节已把这一项枚举为可数的工程产物（RFC、每周发布、VHLO、逐 PR 兼容测试）。
它不体现在 $|E|$ 里，却直接决定采纳率 $p$。

\subsection{本文实证的局限}

表~\ref{tab:threats} 逐条把局限绑定到被削弱的结论上。

\begin{table}[t]
  \caption{本文实证的局限及其对结论的影响。右列说明正文中的处理方式。}
  \label{tab:threats}
  \small
  \begin{tabularx}{\linewidth}{@{}l X X@{}}
    \toprule
    局限 & 证据 & 影响与处理 \\
    \midrule
    语料是上游测试集 & 3,329 个文件中 \texttt{Integration/*} 占 423 个（12.7\%），
      以 GPU 端到端为主 & 投入分布类结论限定为结构性描述，不作使用率推断 \\
    目标端分布不均 & 入矩阵文件的目标家族：通用 CPU 101、AMD 65、通用 GPU 57、Intel 33、
      NVIDIA 21、C/C++ 19、Arm 3，合计 299；另有 72 个落在中间层目标
      & "收敛在 \texttt{llvm}" 表述为\textbf{上游语料中}的收敛；GPU 厂商合计 119，被过度代表 \\
    $|E|$ 是接口数量的代理 & 76 条方言对；一条 \texttt{Conversion} 目录可能包含多个 pass
      & 明写等权假设；结论是结构性的，不是绝对成本 \\
    口径偏差 & \texttt{builtin} 因无命名空间被显著低估；81 个转换文件因一端方言不定被排除
      & 明写 $|E|=76$ 是\textbf{下界}，真实门槛只会更不利 \\
    层划分主观 & 主划分 15 个非零层对，替代划分 20 个；最深层占比 79.2\% $\rightarrow$ 36.6\%
      & 已做稳健性检验；正文只用划分不变量 \\
    只用一个仓库 & 未纳入 \texttt{torch-mlir}、StableHLO、\texttt{onnx-mlir} 的语料
      & 栈顶碎片化只能间接推断，列为后续工作 \\
    参数来源不对称 & $M=11$ 由数据自动实测；$N=8$ 为生态假设；$p$ 只做敏感性
      & 明写三者来源差异；结论不依赖 $p$ 的具体取值 \\
    \bottomrule
  \end{tabularx}
\end{table}

其中两条值得额外说明。

\textbf{$|E|$ 是下界，而这反而加强了结论的方向。}81 个转换文件因一端方言不定被排除，
因此真实方言对数只会比 76 更多；门槛随 $|E|$ 上升，所以工程级门槛只可能更高。
方向上对本文有利，但数值上必须如实标注为下界。

\textbf{"落到最深层"的百分比不可单独引用。}它是 79.2\%（236/298），但这个数字依赖"把硬件目标层
划成一层"这一选择；换成两层划分后掉到 36.6\%。正文因此只用方言级的量（104 个文件、19 个来源
汇入 \texttt{llvm}），不用层级的百分比。

\subsection{结论强度分级}

\textbf{硬结论}（方言级，与层划分无关，可以直接断言）：$|E|=76$（下界）；\texttt{llvm} 的入度为
104 个文件、来自 19 个方言；最深层内部 114 个文件 / 8 条边（68 个厂商内 + 46 个出口转向）；
导入层投入占比 4.5\%。

\textbf{软结论}（依赖参数或划分，必须带条件说）：工程级门槛 $N^{*}\approx 15.5$
（$M=11$、$p=0.6$）；架构级门槛 $N^{*}\approx 4.6$；采纳率敏感性约 2.9 倍。
引用时都需带上 $M$ 与 $p$ 的取值。

\textbf{不可断言}：任何性能、编译时间或使用率数字（本文没有测量）；"MLIR 解决了 $N\times M$"；
单独引用 79.2\%。

\subsection{本文的自我定位}

本文能支持的最强主张是：\textbf{MLIR 把适配成本的组织方式从成对定制改为分层共享，
并在上游语料中呈现出以 \texttt{llvm} 为端点的收敛；而这一转变的工程门槛取决于谁支付 $|E|$}。
本文不能支持的是：这一转变在具体硬件上更快，或已被多数框架采纳。
